\chapter{Discussion}
\label{ch:discussion}

\section{Interpretation of Results}
\label{sec:interpretation}

\subsection{Activation-Level Entity Encoding as Primary Discovery}
Our findings regarding how channels encode multiple entities at different activation strengths reveal a sophisticated information compression mechanism. Rather than requiring separate channels for each entity type, the model has learned to use activation magnitude as an additional dimension for encoding information. This discovery—visible in parallel rollout experiments showing consistent offsets between entity trajectories (correlation 0.931 in conv4a)—represents our most significant finding.

The ability to completely redirect agent behavior through uniform activation offsets (93\% blue key collection shifting to 94\% green key at +4.8 offset) demonstrates this encoding is causally central to network function. This suggests that interventions on neural networks need to consider not just which channels to modify, but also the precise activation values to use.

This activation-level encoding helps explain why the model can maintain full functionality with so few active channels—the network has effectively learned to multiplex information within individual channels. Due to the similarity of the task of tracking each key, it proves efficient to repurpose channels already capable of entity detection, having them target each entity sequentially while indicating the specific entity through activation strength.

\subsection{The Dead Channel Phenomenon in SAEs}
% From Paper 1 Discussion
The finding that a majority of steerable channels in the SAEs were `dead channels' (i.e., channels that exhibit no activation under normal circumstances) is a significant observation. This seemingly counterintuitive result has a clear mechanistic explanation rooted in the SAE architecture used (1x1 convolutional layers). In such an architecture, there is a direct one-to-one relationship between hidden channels in the SAE and their corresponding decoder weights. When a hidden channel dies (ceases to activate during normal forward passes), the gradient flow to its associated decoder weights stops. This effectively freezes those decoder weights at the values they had achieved before the channel became inactive.
Crucially, even though these channels do not activate during normal operation, their frozen decoder weights can still respond to artificial activation during intervention experiments. This explains why steering the agent is possible using channels that are otherwise silent – the learned patterns in their decoder weights are preserved and can be externally triggered. This implies that dead channels should not be automatically dismissed as unimportant in SAE analysis; they may preserve earlier learned representations that offer valuable insights into the model's learning trajectory and latent capabilities.

\subsection{Polysemanticity and Goal Representation in the Base Model}
% From Paper 2 Discussion
While prior work demonstrated interpretable channels in simpler environments (e.g.,\cite{mini_2023_understanding} in Procgen Maze), agents trained in more complex, multi-objective environments like Procgen Heist tend to develop highly polysemantic representations in their base networks. Most early-layer channels, and many in later convolutional layers, activate across multiple, diverse entities, making it challenging to assign singular, clear functional roles. Despite this polysemanticity, steering vectors proved to be a useful tool for isolating goal-related activity. Measuring the sparsity of these vectors allowed for the identification of channels in later layers that reliably track specific goals, even when these goals appear in varied spatial configurations.
These results suggest different layers of abstraction within the model: lower layers appear to capture general visual features and multiple entity types, while later layers consolidate this information into more actionable, goal-directed signals. The methods used offer a diagnostic approach to probing this transition and identifying functionally important parts of the network, even amidst widespread polysemanticity.

\subsection{Solution Multiplicity and Representational Drift}
A striking observation from our experiments was that intervention positions changed completely across different checkpoints even after convergence. Despite maintaining similar performance, the network continued exploring different encoding schemes within the solution space. This is largely due to the fact that we trained the model with an endless variety of maze variations rather than a fixed number, meaning it continued to adapt and change over time. This configuration was also the only one that led to successful training of our smaller model.

This variety of solutions suggests that amplitude-based multiplexing is not a unique solution but rather one of many equivalent representations the network can adopt, and a pattern which is likely to emerge over the course of training. The continuous drift between encoding schemes post-convergence indicates the network exists on a `valley floor' of equally viable solutions, constantly reorganizing its internal representations while preserving behavioral performance. 

We find a notable checkpoint at update step 30001 where no steerable channels at all are found for any entity, though we only found 1 example of this in the 60 checkpoints we tested, suggesting it might be less optimal than those with the stronger intervention configuration.

This has meaningful implications for approaching understanding the inner workings of models. The specific channels and activation ranges we identify for steering are not fixed properties of the task but rather snapshots of a dynamic system exploring equivalent representations as possible solutions. This representational drift may explain some of the difficulty in creating robust interpretability tools and suggests that interventions may need to be continuously recalibrated as networks evolve, though the extent to which this phenomena occurs beyond RL is unknown. 


\subsection{Activation-Level Entity Encoding}
Our findings regarding how the intervention spans encode multiple entities at different activation strengths reveal a sophisticated information compression mechanism. Rather than requiring separate channels for each entity type, the model has learned to use activation magnitude as an additional dimension for encoding information. This suggests that interventions on neural networks may need to consider not just which channels to modify, but also the precise activation values to use.

This activation-level encoding also helps explain why the model can maintain full functionality with so few active channels, the network has effectively learned to multiplex information within individual channels. 

An explanation for why this happens might be that due to the similarity of the task of tracking each key, it proves the most efficient solution to have re-purpose channels already capable of entity detection in general, and to have them simply target each one sequentially, while indicating the nature of the given entity by varying the activation strength.

\subsection{Role specialization}

A surprising finding was that some channels seem able to unilaterally alter where the agent will navigate to, many other channels were more reliably able to navigate the agent when the other channels were ablated, but were completely unable to individually steer the agent, as evidenced by examining \ref{fig:steering-dot-plot} and \ref{fig:ablation-heatmap}. This indicates that the specific function of indicating that the agent `must go to this spot' were exclusive to the steering channels, and the channels that were more successful in navigating had a role specifically in identifying how to get to a position, that in the absence of any other channels providing contrasting signals would lead the agent to go to that position. 

This idea of specialization and reliance between channels is reinforced by the fact that ablating both the navigation channels and the intervention spans produced better collection rates than ablating the best navigation channels while leaving the intervention spans untouched. The fact that this occurred 52\% of cases with n=400 trials is evidence that this effect is robust and not just noise.


\section{Theoretical Implications}
\label{sec:theoretical}
% How your results contribute to theory
% New insights about RL interpretability
The findings from this research contribute to the theoretical understanding of mechanistic interpretability in deep reinforcement learning in several ways. The demonstration of activation-level encoding suggests that neural networks may employ more nuanced information encoding strategies than simple feature presence/absence, utilizing the analog nature of activations as an additional information dimension. This challenges some simpler models of feature representation and highlights the potential for information multiplexing as a core principle of neural computation, especially when under pressure for representational efficiency.

The dead channel phenomenon in SAEs, and its mechanistic explanation, provides a concrete example of how network components that appear non-functional during standard inference can retain latent capabilities and learned knowledge. This has implications for how we assess model capacity and robustness, and for understanding the learning dynamics that lead to such states.

One observation from carrying out this work is that uncovering the role of activation magnitude across entire layers as a selection mechanism was difficult to uncover using typical mechanistic interpretability tools such as probes and SAEs. The most important discovery regarding how the network functioned (activation level entity encoding) turned out to be explicable purely by analysing the mean activations of different channels. It is possible that there may be further interesting discoveries to be made using similar approaches, particularly in networks that would likely benefit from recycling certain functionality for similar tasks that require clear delineation.

For RL interpretability specifically, this work underscores that techniques successful in one domain (e.g., LLMs) or simpler environments may need adaptation or re-evaluation in more complex, interactive settings. The interplay between learned environmental structure (like sequential objectives in Heist) and emergent network representations is a key area for future theoretical development.

\section{Practical Applications}\label{sec:applications}
% Real-world applications of your methods
% Potential use cases in different domains
While this research is primarily foundational, aiming to understand the internal workings of RL agents, the insights gained could have several long-term practical applications. A deeper understanding of how agents represent goals and make decisions can inform the development of more robust and reliable RL systems. For instance, identifying critical control channels or specific activation patterns linked to desired (or undesired) behaviors could lead to methods for:
\begin{itemize}
    \item \textbf{Enhanced Safety and Control:} Modifying or monitoring these identified neural correlates could provide mechanisms to prevent catastrophic failures or to guide agent behavior in critical situations.
    \item \textbf{Improved Debugging and Diagnostics:} When an RL agent fails or behaves unexpectedly, the ability to inspect specific features or channels (especially those disentangled by SAEs) could significantly speed up the debugging process.
    \item \textbf{More Sample-Efficient Learning:} Understanding what representations are crucial for a task might eventually lead to techniques that encourage or bootstrap the learning of these representations, potentially improving sample efficiency.
    \item \textbf{Agent Alignment:} Identifying how an agent represents complex or potentially conflicting goals is a step towards ensuring that agents behave in accordance with human intentions, a key challenge in AI alignment.
\end{itemize}
The techniques themselves, such as targeted activation steering or the analysis of SAE latents, could become part of a standard toolkit for RL practitioners seeking to verify and validate their models beyond simple performance metrics.

\section{Limitations of the Study}\label{sec:limitations}
% Honest assessment of limitations
% Factors that might affect validity or generalizability
This research, while providing several insights, is subject to certain limitations:
\begin{itemize}
    \item \textbf{Specificity to Architecture and Environment:} The findings are based on specific CNN architectures trained on the Procgen Heist environment. While efforts were made to use simplified, interpretable models and to validate on standard architectures (Paper 2), the extent to which these detailed mechanistic findings generalize to vastly different architectures (e.g., Transformers, RNNs) or other complex RL environments remains an open question.
    \item \textbf{Entangled Activations:} As noted in Paper 2, many goal-related activations in the base model remain entangled across channels. While SAEs aim to disentangle these, the disentanglement may not be perfect, and some interpretations might still be confounded by residual polysemanticity or distributed representations.
    \item \textbf{Off-Distribution Behavior:} Intervention experiments, by their nature, can push the network into off-distribution activation states. While this is useful for causal testing, the agent's behavior under such strong artificial interventions might not fully reflect its operational characteristics under normal, on-distribution conditions.
    \item \textbf{Incomplete Mapping of Channel Functions:} While specific channels and their roles were identified, a complete, precise map of all channels and their exact contributions to various sub-tasks or representations was not achieved. The complexity of the network, even simplified ones, makes such exhaustive mapping a significant challenge.
    \item \textbf{SAE Interpretation Nuances:} The interpretation of SAE latents themselves can be complex. It is not always clear if the features learned by an SAE perfectly mirror the computations of the base model or if they represent a re-factorization that is useful but distinct. The stability of SAE training and its impact on the consistency of learned features across different training runs or minor hyperparameter changes is also a consideration (Paper 1).
    \item \textbf{Focus on Convolutional Layers:} The primary focus of intervention and SAE application was on the convolutional layers. While these are crucial for visual processing and early planning, a full understanding would also require similar deep dives into the fully connected layers and the policy/value heads.
\end{itemize}

Our work is currently limited by the fact that we test only on a single environment, as there are not many environments that have similar qualities such as the extreme generalization encouraged by Procgen, and the sequential goals that Heist contains. To address this, future work could attempt this with MiniGrid environments, though a custom configuration that can establish similar diversity may be needed.

It is also unclear to what extent the intervention channel multiplexing phenomena generalises to other architectures, and we plan to test this in future work.

We also fail to provide an explanation of the specific mechanism played by the intervention channels. We plan to address this directly through a thorough application of probes that are trained to detect which entity is currently the next entity in order of pursuit by the network. Early work shows that the later convolutional layers do this very successfully, as well as the first fully connected layer. However, to truly understand the role of the intervention spans and navigation channels, it may be necessary to train probes on every channel and every row in the fully connected layer to determine their function. 

Our approach of simply intervening in every channel with 400 different activation strengths is also an approach that would need to be approached more carefully, due to cost constraints, in larger models. 

\section{Future Research Directions}\label{sec:future}
% Suggestions for further research
% Unsolved problems and next steps
The findings and limitations of this work point to several promising avenues for future research:

% From Paper 1 Future Work
\begin{itemize}
    \item \textbf{Tracking Non-Current Objectives:} Further exploration is needed into how RL agents internally represent and track entities or sub-goals that are not the immediate next target in a sequence, and whether interventions can reliably switch the agent's focus to these deferred objectives.
    \item \textbf{Application to Different Architectures:} Extending this line of inquiry to more complex agent architectures, such as decision transformers or RNN-style networks, would be valuable. These architectures have different mechanisms for handling temporal information and memory, and training probes or performing targeted interventions in their residual streams or hidden states could yield new insights.
    \item \textbf{Deeper Investigation of Activation Strength Encoding:} The phenomenon of activation strength encoding entity identity within SAE channels warrants further investigation. Understanding the precise mechanisms and robustness of this encoding could be highly beneficial.
    \item \textbf{Exploring Other Environment Entities:} This work focused primarily on keys and the gem. Future studies could investigate if other visual cues (e.g., the image of collected keys in a UI element) trigger steerable behaviors related to corresponding objectives (e.g., locks).
    \item \textbf{Advanced Attribution Techniques:} Applying more advanced techniques like attribution-based parameter decomposition~\citep{braun_2025_interpretability} to networks like the one studied here could yield valuable and potentially generalizable insights into how parameters contribute to function.
    \item \textbf{SAEs and Base Model Correspondence:} Investigating whether the functionalities highlighted or preserved by SAEs (such as entity encoding via activation strength) truly reflect similar underlying mechanisms in the base model, or are emergent properties of the SAE compression process, is a critical follow-up question.
\end{itemize}

% From Paper 2 Future Work / Discussion points framed as future work
\begin{itemize}
    \item \textbf{Robust Isolation of Goal-Related Circuits:} While steering vectors and SAEs help, further development of tools like causal tracing or more advanced sparse coding methods could be effective in more robustly isolating complete goal-related circuits, especially in the presence of significant polysemanticity.
    \item \textbf{Impact of Environment Design on Learned Representations:} Investigating how variations in environment design (e.g., training models with no fixed ordering of objectives in the maze) affect the development of generalizable circuits for goal pursuit would be an interesting direction.
    \item \textbf{Synthesizing Distributed Information:} Developing tools or models capable of synthesizing information from across different parts of the network is crucial, as the factors contributing to a specific outcome are often distributed. SAEs are one step in this direction, but further methods are needed.
\end{itemize}

\section{Reflections on the Research Process}\label{sec:reflections}
% Personal insights gained
% Challenges encountered and overcome
(This section is for your personal reflections on the research journey, challenges faced, and lessons learned. This content should be written by you.)